\documentclass[twocolumn]{aastex631}
\begin{document}
\title{A residual reionization ionizing-photon-budget shortfall for star-forming galaxies at z\~{}8 (highC systematic corner)}
\author{NebulaMind Lab (autonomous pipeline)}
\affiliation{NebulaMind Lab, \url{https://lab.nebulamind.net}}
\begin{abstract}
Reionization ionizing-photon-budget reconciliation at z\~{}8: star-forming galaxies require f\_esc=0.220 (+0.206/-0.108) to close the budget (Madau-Dickinson SFRD, log xi\_ion=25.5+/-0.15, clumping C=2-5, JWST-SFRD tail), versus indirect-proxy-inferred f\_esc=0.062 (+0.108/-0.039) from LzLCS O32/beta calibrations. Median delta(required-inferred)=+0.140 dex-frac (16-84\%: +0.005 to +0.347); 85\% of the systematic MC shows a shortfall. Conclusion: a genuine SHORTFALL remains, and the sign holds under both O32 and beta calibrations. The result is bounded by the xi\_ion x clumping x proxy-calibration systematic, not by statistics. Generated autonomously from public data (jwst) via the NebulaMind Lab runner: a bounded, reproducible, descriptive study.
\end{abstract}
\section{Introduction}
Recent studies have highlighted a potential crisis in understanding the reionization process, with concerns over whether star-forming galaxies can produce enough ionizing photons to drive this cosmic event [Muñoz2024]. This issue is further complicated by uncertainties in the escape fraction of these photons, as well as the role of other factors such as clumping and the ionizing efficiency of galaxies [Davies2021], [Park2022]. To address this problem, we need to reconcile the reionization photon budget using existing literature values for key parameters.
\section{Data and method}
In our analysis, we rely on published values for the cosmic star formation rate density (SFRD) from Madau \& Dickinson (2014), as well as calibrations for the ionizing efficiency (xi\_ion) and escape fraction (f\_esc) proxies. Specifically, we use the LzLCS O32/beta calibrations from Chisholm+22 and Flury+22, and the Simmonds+24 calibration. We do not employ any new observational data or survey catalog information in this study.
\section{Result}
Our reconciliation of the reionization ionizing-photon budget at z\~{}8 indicates that star-forming galaxies require a higher escape fraction (f\_esc=0.220 +0.206/-0.108) to account for the necessary photons, compared to the indirect-proxy-inferred value (f\_esc=0.062 +0.108/-0.039) derived from LzLCS O32/beta calibrations. This results in a median shortfall of +0.140 dex-frac, with 85\% of systematic Monte Carlo simulations showing a deficit.
\begin{figure}\centering\includegraphics[width=\columnwidth]{result.png}\caption{A residual reionization ionizing-photon-budget shortfall for star-forming galaxies at z\~{}8 (highC systematic corner)}\end{figure}
\section{Caveats}
It is essential to acknowledge the limitations of our approach. Our analysis relies on literature values for key parameters, which may introduce uncertainties due to variations in assumptions and methodologies across different studies. Additionally, our use of automated single-selection and uncalibrated measurements may not capture the full complexity of the reionization process. Furthermore, we have not considered potential systematic errors in the calibrations used or the impact of other factors such as dust attenuation and galaxy evolution on the ionizing photon budget. These limitations highlight the need for further research and more robust observational data to refine our understanding of reionization.
\section*{References}
\footnotesize
[Muoz2024] Muñoz et al. (2024). Reionization after JWST: a photon budget crisis?. bibcode 2024MNRAS.535L..37M \par [Davies2021] Davies et al. (2021). The Predicament of Absorption-dominated Reionization: Increased Demands on Ionizing Sources. bibcode 2021ApJ...918L..35D \par [Park2022] Park et al. (2022). Calibrating excursion set reionization models to approximately conserve ionizing photons. bibcode 2022MNRAS.517..192P \par [Duncan2015] Duncan et al. (2015). Powering reionization: assessing the galaxy ionizing photon budget at z < 10. bibcode 2015MNRAS.451.2030D \par [Madau2017] Madau (2017). Cosmic Reionization after Planck and before JWST: An Analytic Approach. bibcode 2017ApJ...851...50M

\end{document}
